\chapter{A soliton, and the boundary it shares}

\textbf{Purpose:} Stop squaring the projection, find the object that squaring hid, and take it apart
until every number in it is accounted for. \textbf{Scope:}
\texttt{src/\allowbreak{}bench/\allowbreak{}bench\_\allowbreak{}sac.\allowbreak{}cu}, the
\texttt{signed} and \texttt{cells} arms.

\section{What squaring destroyed}

Every projection in the preceding chapters sums \(z^{2}\) along an axis, and \(z^{2}\) cannot
separate a rate of zero from a rate of one: both give \(z = \mp\sqrt{n}\) and the same square. The
sum is blind to sign. A cell that changes nothing and a cell that flips its output \emph{every
single time} land on the same value and are drawn at the same height in the same color.

The obvious repair is worse. Summing signed \(z\) makes a row that is half never-flipping and half
always-flipping come to zero, the same total a fully mixed row gives, and hides the same
distinction one level down.

So the three states are counted and not summed. A count of zero makes the cell \emph{never}; a
count of every trial makes it \emph{always}; any count between the two makes it \emph{mixed}. Those
three are exclusive and cover every case, and nothing cancels.

\section{Something that does not move while everything else does}

\begin{center}
\begin{tabular}{rrrr}
\toprule
round & never & always & mixed \\
\midrule
6 & 99{,}051 & \textbf{312} & 31{,}709 \\
10 & 66{,}294 & \textbf{312} & 64{,}466 \\
16 & 17{,}126 & \textbf{312} & 113{,}634 \\
17 & 9{,}840 & 248 & 120{,}984 \\
21 & 0 & 0 & 131{,}072 \\
\bottomrule
\end{tabular}
\end{center}

The \emph{always} count holds at exactly \(312\) for eleven consecutive rounds while \emph{never}
falls by a factor of \(5.8\) and \emph{mixed} climbs by eighty-two thousand around it.

It is not standing still in the field. It is travelling.

\begin{center}
\begin{tabular}{rlrrl}
\toprule
round & bits & width & cells & words \\
\midrule
6 & 57--191 & 135 & 312 & 1--5 \\
8 & 121--255 & 135 & 312 & 3--7 \\
10 & 185--319 & 135 & 312 & 5--9 \\
12 & 249--383 & 135 & 312 & 7--11 \\
14 & 313--447 & 135 & 312 & 9--13 \\
16 & 377--511 & 135 & 312 & 11--15 \\
\bottomrule
\end{tabular}
\end{center}

\textbf{A rigid object, one hundred and thirty-five bits wide and three hundred and twelve cells in
area, moving at exactly thirty-two bits per round without changing form.} That is a soliton: a shape
that propagates without dispersing. It is stationary in the co-moving frame, and so no reading
that averaged over rounds could see it, and it sits at the same slope of thirty-two that the
slant-stack found statistically at \(7.18\) sigmas, except that this is exact, with no variance
across eleven rounds.

\section{Taking it apart}

The projection sums over the output axis, so it cannot say which output bits these are. Naming them
accounts for every number.

\begin{center}
\begin{tabular}{lrl}
\toprule
part & count & mechanism \\
\midrule
aligned, residue 0 & 256 & integer addition \\
skewed, residue 25 & 56 & \(\Sigma_1\)'s \(\mathrm{ROTR}25\), alone \\
\bottomrule
\end{tabular}
\end{center}

The aligned half is arithmetic. Flipping bit \(k\) of an addend always flips bit \(k\) of the sum,
because carries only run upward, so bit \(k\) of a message word drives bit \(k\) of \(T_1\) and
thence of \(a\) and of \(e\), two cells per bit, over four words of thirty-two bits:
\(4 \times 32 \times 2 = 256\).

The word breakdown closes too. The newest word contributes \(64\) and is purely aligned; the three
behind it contribute \(78\) each, being \(64\) aligned and \(14\) skewed; the oldest contributes
\(14\), skewed only, its aligned cells already mixed away. Four times sixty-four is two hundred and
fifty-six, four times fourteen is fifty-six, and the certain cells spread perfectly evenly over the
eight state words, thirty-nine each, at every depth.

\subsection{Where the explanation fails}

Every one of the fifty-six skewed cells sits at residue twenty-five. None at \(\mathrm{ROTR}6\),
none at \(\mathrm{ROTR}11\), none at any of \(\Sigma_0\)'s amounts, at rounds eight, ten and twelve
alike.

The natural mechanism does not survive that. If carries destroy certainty at every position above
the lowest flipped one, then the lowest of \(\Sigma_1\)'s three images should be the survivor, and
at source bit twenty that is \(k - 11 = 9\), below \(k - 25 = 27\), so residue eleven should appear.
It never does.

\textbf{The aligned half is understood and the skewed half is an exact measurement with a wrong
explanation.} That is worth more than a vague open question, because it names the observation the
next model has to reproduce.

\section{The boundary this shares with the published work}

A soliton of probability-one paths is what a local-collision construction is built out of, so this
object has neighbors in the published work. An earlier draft of this section claimed a numerical
coincidence that does not exist, and the correction is more interesting than the claim was.

\subsection{A retracted claim}

That draft read: \emph{Sanadhya and Sarkar construct deterministic twenty-one step collisions, and
the signed projection says probability-one propagation ends at round twenty-one; two opposite
methods on one number.} It was written from a search-engine summary and not from the paper.

\textbf{The deterministic limit is twenty-two, not twenty-one.} The same two authors passed
twenty-one within months of the paper cited, and record it themselves: ``Deterministic constructions
of up to 22-round SHA-2 collisions are described using the SS local collision.'' The match was a
coincidence resting on a wrong number, and it is retracted.

\subsection{The boundary that is real}

Their limit is not arbitrary and its mechanism is a boundary this work found independently. A local
collision spans nine steps, and a deterministic construction requires it to sit entirely inside the
sixteen free message words, since those can be solved for directly while every later word is a
function of them. So it cannot begin after step seven (``obtaining deterministic collisions up to
22 rounds did not require the (single) local collision to extend beyond Step 15''), and beyond
that some expanded word necessarily carries a difference and determinism ends.

That is the message expansion switching on, which the schedule ablation in the previous chapter
locates at round seventeen, bit for bit. \textbf{The published deterministic ceiling and the knee
measured here are the same structural fact}, approached from construction and from measurement.
Round twenty-one, where the last certain cell disappears, is a different and later boundary.

\subsection{What the sources actually say}

\textbf{On the most significant bit.} Sanadhya and Sarkar state that ``since the XOR and additive
differences coincide for the most significant bit, to achieve higher probability, it is advantageous
to ensure that many bits in the differential path are MSBs'', because a one-bit difference survives
addition ``with probability one if the differing bit is the most significant bit, else with
probability half''.

That is \emph{adjacent to} the aligned half of this object and not the same statement, and an
earlier draft ran them together. Their claim is that a difference stays weight one only at the MSB.
The claim here is that bit \(k\) of an addend always flips bit \(k\) of the sum, which holds at every
\(k\); what the MSB adds is that no carry leaves it. Both follow from the carry running one way, but
they are different facts and only the second is what the \(256\) aligned cells count.

\textbf{On the message expansion.} ``The message expansion of SHA-2 does not play any role in first
16 steps'', stated near-verbatim, and the local collisions of that work span nine steps.

\textbf{On the avalanche measurements.} Vaughn and Borowczak test all \(127\) non-empty combinations
of the seven sub-functions, and name \(\Sigma_1\), integer addition, Choose and the message scheduler
as those contributing to the avalanche criterion at the earliest rounds throughout. That is the
\(T_1\) asymmetry of the operations chapter, found by ablation and not by spectrum, and it is
their result before it was this one's.

Their round figures want quoting carefully. Full SHA-256 exhibits the criterion at round twenty-
three, \emph{fails it between rounds fifty-four and fifty-seven}, and exhibits it again at
fifty-eight. Twenty-three is the floor across combinations, not the full function's plateau.

\textbf{And this work disagrees with them about the dip.} The refutations chapter records the
fifty-four to fifty-seven failure as not reproducing at roughly four times their sample. That
disagreement is left standing and not reconciled: their threshold is a Bonferroni bound at
ninety-nine per cent over \(131{,}072\) simultaneous tests, and the dip sits close enough to it that
sample size decides the answer.
